\documentclass[11pt,a4paper]{article}
\usepackage[utf8]{inputenc}
\usepackage[T1]{fontenc}
\usepackage[spanish]{babel}
\usepackage{amsmath,amssymb,mathtools}
\usepackage{geometry}
\usepackage{enumitem}
\usepackage{lmodern}
\usepackage{microtype}
\geometry{margin=2.4cm}
\setlength{\parindent}{0pt}
\setlength{\parskip}{0.5em}
\newcommand{\R}{\mathbb{R}}
\newcommand{\Span}{\operatorname{L}}
\newcommand{\Ima}{\operatorname{Im}}
\newcommand{\sen}{\operatorname{sen}}
\begin{document}

\begin{center}
{\LARGE\bfseries Álgebra Lineal}\\[0.4em]
{\Large Tema 4: El Espacio Afín}
\end{center}

\tableofcontents
\newpage

\section{Espacio afín}

Sea $V$ un espacio vectorial. Un conjunto no vacío $\mathcal A$ es un espacio afín sobre $V$ si a cada par $P,Q\in\mathcal A$ le corresponde un único vector $\overrightarrow{PQ}\in V$ y:

\[
\overrightarrow{PQ}+\overrightarrow{QR}=\overrightarrow{PR}.
\]

Equivalentemente, existe una operación:

\[
\mathcal A\times V\longrightarrow\mathcal A,
\qquad
(P,v)\longmapsto P+v,
\]

tal que para cada $P,Q\in\mathcal A$ existe un único $v\in V$ con $P+v=Q$ y:

\[
(P+u)+v=P+(u+v).
\]

Los elementos de $\mathcal A$ se denominan puntos y $V$ es su espacio de dirección.

\[
\boxed{\dim\mathcal A=\dim V}
\]

\subsection{Propiedades}

\[
P+v=P\iff v=0
\]

\[
\overrightarrow{QP}=-\overrightarrow{PQ}
\]

\[
\boxed{\overrightarrow{PQ}+\overrightarrow{QR}=\overrightarrow{PR}}
\]

La última igualdad es la relación de Chasles.

Si:

\[
\overrightarrow{PQ}=\overrightarrow{P'Q'},
\]

entonces:

\[
\boxed{\overrightarrow{PP'}=\overrightarrow{QQ'}}.
\]

\section{Espacio afín métrico}

Si $V$ es un espacio vectorial euclídeo, $\mathcal A$ es un espacio afín métrico.

La distancia entre dos puntos es:

\[
\boxed{
d(P,Q)=\left\|\overrightarrow{PQ}\right\|
}
\]

y cumple:

\[
d(P,Q)=0\iff P=Q
\]

\[
d(P,Q)=d(Q,P)
\]

\[
d(P,R)\leq d(P,Q)+d(Q,R).
\]

En $\R^n$:

\[
\boxed{
d(P,Q)=
\sqrt{\sum_{i=1}^{n}(b_i-a_i)^2}
}
\]

si $P=(a_1,\ldots,a_n)$ y $Q=(b_1,\ldots,b_n)$.

El ángulo entre dos vectores no nulos viene dado por:

\[
\boxed{
\cos\alpha=
\frac{u\cdot v}{\|u\|\|v\|}
}
\]

\section{Sistemas de referencia. Coordenadas}

Un sistema de referencia cartesiano de un espacio afín de dimensión $n$ es:

\[
\boxed{
R=\{O,B\}
}
\]

donde $O$ es un punto y:

\[
B=\{u_1,\ldots,u_n\}
\]

es una base del espacio de dirección.

Si:

\[
\overrightarrow{OP}=x_1u_1+\cdots+x_nu_n,
\]

entonces:

\[
\boxed{
[P]_R=(x_1,\ldots,x_n).
}
\]

Dados:

\[
[P]_R=(a_1,\ldots,a_n),
\qquad
[Q]_R=(b_1,\ldots,b_n),
\]

se tiene:

\[
\boxed{
[\overrightarrow{PQ}]_B=
(b_1-a_1,\ldots,b_n-a_n).
}
\]

Si $B$ es ortonormal, $R$ es un sistema de referencia rectangular.

\section{Cambio de sistema de referencia}

Sean:

\[
R=\{O,B\},
\qquad
R'=\{O',B'\}.
\]

Si:

\[
[O']_R=b
\]

y $M_{B'B}$ tiene por columnas los vectores de $B'$ expresados en $B$, para un punto $P$:

\[
[P]_R=x,
\qquad
[P]_{R'}=y,
\]

se cumple:

\[
\overrightarrow{OP}
=
\overrightarrow{OO'}
+
\overrightarrow{O'P}.
\]

Por tanto:

\[
\boxed{
x=b+M_{B'B}y
}
\]

y:

\[
\boxed{
y=M_{B'B}^{-1}(x-b).
}
\]

En coordenadas homogéneas:

\[
\boxed{
\begin{pmatrix}1\\x\end{pmatrix}
=
\begin{pmatrix}
1&0\\
b&M_{B'B}
\end{pmatrix}
\begin{pmatrix}1\\y\end{pmatrix}.
}
\]

\section{Variedades afines}

Sea $P\in\mathcal A$ y $S\leq V$. La variedad afín que pasa por $P$ y tiene espacio de dirección $S$ es:

\[
\boxed{
F=P+S.
}
\]

Si:

\[
S=\Span\{v_1,\ldots,v_k\},
\]

entonces:

\[
\boxed{
F=
P+\Span\{v_1,\ldots,v_k\}
}
\]

y:

\[
\boxed{\dim F=k.}
\]

\subsection{Ecuaciones paramétricas}

\[
\boxed{
X=P+\lambda_1v_1+\cdots+\lambda_kv_k.
}
\]

\subsection{Ecuaciones cartesianas}

Una variedad afín es el conjunto de soluciones de un sistema compatible:

\[
AX=b.
\]

Si $P$ es una solución particular:

\[
AP=b,
\]

entonces:

\[
A(X-P)=0.
\]

Así:

\[
\boxed{
F=P+\Ker A.
}
\]

\section{Independencia afín}

Los puntos:

\[
P_0,P_1,\ldots,P_k
\]

son afínmente independientes si:

\[
\boxed{
\overrightarrow{P_0P_1},
\ldots,
\overrightarrow{P_0P_k}
\text{ son linealmente independientes}.
}
\]

Un conjunto de $n+1$ puntos afínmente independientes determina un sistema de referencia afín de dimensión $n$.

\section{Variedad afín generada por puntos}

Dado:

\[
\mathcal P=\{P_0,P_1,\ldots,P_k\},
\]

la menor variedad afín que los contiene es:

\[
\boxed{
L(\mathcal P)
=
P_0+
\Span\{
\overrightarrow{P_0P_1},
\ldots,
\overrightarrow{P_0P_k}
\}.
}
\]

En $\R^3$, una recta se determina por dos puntos distintos y un plano por tres puntos no alineados.

\section{Intersección y suma de variedades afines}

Sean:

\[
F_1=P_1+S_1,
\qquad
F_2=P_2+S_2.
\]

Si:

\[
F_1\cap F_2\neq\emptyset
\]

y $P\in F_1\cap F_2$, entonces:

\[
\boxed{
F_1\cap F_2=P+(S_1\cap S_2)
}
\]

y:

\[
\boxed{
F_1+F_2=P+(S_1+S_2).
}
\]

Además:

\[
\boxed{
\dim F_1+\dim F_2
=
\dim(F_1\cap F_2)+\dim(F_1+F_2).
}
\]

Si la intersección es vacía, la variedad generada por ambas tiene como espacio de dirección:

\[
\boxed{
\Span\{
S_1,S_2,\overrightarrow{P_1P_2}
\}.
}
\]

\section{Aplicaciones afines}

Sean $\mathcal A$ y $\mathcal A'$ espacios afines, con espacios de dirección $V$ y $V'$.

Una aplicación:

\[
f:\mathcal A\longrightarrow\mathcal A'
\]

es afín si existe una aplicación lineal:

\[
\widetilde f:V\longrightarrow V'
\]

tal que:

\[
\boxed{
\widetilde f(\overrightarrow{PQ})
=
\overrightarrow{f(P)f(Q)}
}
\]

para todos $P,Q\in\mathcal A$.

Fijando un punto $O$:

\[
\boxed{
f(X)=f(O)+\widetilde f(\overrightarrow{OX}).
}
\]

La composición de aplicaciones afines es afín.

Además:

\[
\boxed{
f\text{ es inyectiva}
\iff
\widetilde f\text{ es inyectiva}
}
\]

y:

\[
\boxed{
f\text{ es suprayectiva}
\iff
\widetilde f\text{ es suprayectiva}.
}
\]

\section{Expresión matricial de una aplicación afín}

Sean:

\[
R=\{O,B\},
\qquad
R'=\{O',B'\}.
\]

Si:

\[
[f(O)]_{R'}=C
\]

y:

\[
A=M_{B,B'}(\widetilde f),
\]

entonces:

\[
\boxed{
Y=C+AX.
}
\]

En coordenadas homogéneas:

\[
\boxed{
\begin{pmatrix}
1\\Y
\end{pmatrix}
=
\begin{pmatrix}
1&0\\
C&A
\end{pmatrix}
\begin{pmatrix}
1\\X
\end{pmatrix}.
}
\]

\end{document}
